% TOP_PROBLEM: {"id": 3800006, "title": "Halving lines and k-sets", "category_id": 6, "category": {"id": 6, "name": "geometry", "display_name": "Geometry", "description": "Euclidean and non-Euclidean geometry, geometric structures.", "slug": "geometry", "order_index": 6, "created_at": "2026-07-31T15:26:25.671Z"}, "status": "partially_solved", "problem_number": "AMR-037-0006", "set_id": 13}
% FIRST_POSTED: 2026-10-01
% ENTRY_KIND: statement_only
% UnsolvedMath ID 3800006; source catalogue code AMR-037-0006
% !TeX program = lualatex
% Definition and sources only; this file is not an LLM solution attempt.
\documentclass[11pt,a4paper]{article}
\usepackage[margin=25mm]{geometry}
\usepackage{fontspec}
\setmainfont{lmroman10-regular.otf}[BoldFont=lmroman10-bold.otf,
 ItalicFont=lmroman10-italic.otf,BoldItalicFont=lmroman10-bolditalic.otf]
\usepackage{amsmath,amssymb,mathtools}
\usepackage{unicode-math}
\setmathfont{latinmodern-math.otf}
\AtBeginDocument{\let\setminus\smallsetminus}
\usepackage{xurl,hyperref}
\hypersetup{colorlinks=true,urlcolor=blue}
\setcounter{secnumdepth}{0}
\setlength{\parindent}{0pt}
\setlength{\parskip}{5pt}
\setlength{\emergencystretch}{3em}
\begin{document}
\section*{Halving lines and k-sets}\label{3800006}
{\small UnsolvedMath ID 3800006 · AMR-037-0006 · Geometry · AMR Open Problem Lists}
\subsection{Definitions and mathematical statement}
Let $S\subset\mathbb R^2$ consist of $n$ distinct points, with no three
collinear. A $k$-set is a subset $A\subset S$ of size $k$ for which some
line disjoint from $S$ strictly separates $A$ from $S\setminus A$.
Count distinct subsets, not the infinitely many separating lines. Define
\[
 K_2(n,k)=\max_{|S|=n}\#\{A\subset S:A\text{ is a }k\text{-set}\}.
\]
Find sharp bounds as functions of both $n$ and $1\le k\le n-1$.

For even $n$, a halving line is a line through two points of $S$ with
exactly $(n-2)/2$ of the remaining points in each open half-plane. Determine
the largest number $H_2(n)$ of these lines. The halving-line count is a
closely related extremal question; its objects are not the separating
lines in the preceding definition.

In $\mathbb R^d$, take sets in affine general position, define $k$-sets by
strict separation with a hyperplane, and determine the analogous
$K_d(n,k)$. A halving hyperplane passes through $d$ points and has equally
many remaining points on each side, so this equal-split version requires
$n-d$ even. Count each resulting affine hyperplane once.
\subsection{Short English statement}
Determine how many subsets of a point configuration can be cut off by
a line or hyperplane, with particular emphasis on equal-size splits.
\subsection{Sources}
\begin{itemize}
\item[S1] ULAM.AI and the UnsolvedMath Contributors, supplied problems.json, record 3800006 (AMR-037-0006), AMR Open Problem Lists. Catalogue identity and source transcription; CC BY 4.0.
\url{https://huggingface.co/datasets/ulamai/UnsolvedMath}.
\item[S2] Erickson - Open Problems in Computational Geometry, Halving lines and k-sets. Original source indexed by AMR-037-0006.
\url{https://jeffe.cs.illinois.edu/open/}.
\end{itemize}
\end{document}
